% ===== 原文 PDF pp.47--50（印刷 pp.46--49） =====
% tr-04 的 tr04c 片段。p.47（印刷 p.46）页首是表 3（归属 tr-04）。
% 本文件 tag：3.42（印刷 p.46）；4.1（印刷 p.48）。共 2 条。
% 终点：p.50（印刷 p.49）页末。—— 注意：实测 p.50 上除 §4.2 标题外，还含 §4.2 的全部正文
% （四段）与脚注 34；而 (4.2)(4.3) 在 p.51，属 §4.3（tr-05）。按 STYLE_GUIDE §8「按实际页面写」，
% 本块译全 p.50，并已在回执中报告此边界与 BRIEF 不符。

% [原文 p.46]
\begin{center}
\begin{tabular}{lcc}
\toprule
理论 & 味群 & $\Delta[u]$ \\
\midrule
$H_0$ & $-$ & $6/5$ \\
$H_1$ & $\SU(2)$ & $4/3$ \\
$H_2$ & $\SU(3)$ & $3/2$ \\
$D_4$ & $\SO(8)$ & $2$ \\
$\mathrm{E}_6$ & $\mathrm{E}_6$ & $3$ \\
$\mathrm{E}_7$ & $\mathrm{E}_7$ & $4$ \\
$\mathrm{E}_8$ & $\mathrm{E}_8$ & $6$ \\
\bottomrule
\end{tabular}
\end{center}
\tabtitle{3}{各种秩一 SCFT 中手征主算子 $u$ 的共形维数。理论 $H_0$ 是最简单的 AD 理论，即 $(A_1,A_2)$，其整体对称性是平凡的。这些理论的名字指的是 F 理论中 D3 膜所探测的奇点（对 $H_n$ 型奇点而言即 $A_n$）[171--180]。这些理论的其他细节将在 5.4.2 节给出。}

% [Tutor Note: 表 3 三列依原文为「理论 / 味群 / $\Delta[u]$」，逐格核对：$H_0$ 无味对称性（记作 $-$）、$\Delta=6/5$；$H_1$、$H_2$ 的味群为 $\SU(2)$、$\SU(3)$；$D_4$ 为 $\SO(8)$；$\mathrm{E}_{6,7,8}$ 的味群即其自身。表题中 $H_n$ 型奇点对应 $A_n$，但 $H_0$ 是例外（$H_0=(A_1,A_2)$，味对称性平凡）。]

最典型的例子是 $\cN=4$ SYM（视作带一个伴随超多重态的 $\cN=2$ 理论）以及与四个基本超多重态耦合的 $\cN=2$ 理论。它们在各自的库仑分支上具有相同的 SW 几何 $(\Sigma_u,\lambda)$ [21]，但它们是完全不同的 4d 理论。

% [Tutor Note: 这是 [21] 中的经典反例：两套理论的 $(\Sigma_u,\lambda)$ 完全相同，却是不同的 4d 理论。原因在于 SW 几何只编码库仑分支上的低能有效数据，不含 UV 完备化信息——这正是下一段要引入“质量形变等额外数据”的理由。]

的确，为了从库仑分支 EFT 中提取更多关于 SCFT 的信息，有必要研究 SW 几何的一般形变，例如质量形变。特别地，由于库仑分支由定义 2.1 中所列手征主算子的 VEV 参数化，仅从 EFT 就有可能提取这些受保护算子的谱（正如我们在上一节对秩一情形所回顾的）。耦合与质量形变也自库仑分支 EFT 中显现，并且可以从完全形变后的 SW 曲线读出。

从库仑分支 EFT 中还能提取的另一类 SCFT 数据，是它的共形中心荷与味中心荷。超对称 Ward 恒等式把这些 CFT 可观测量与涉及 $\mathrm{U}(1)\times\SU(2)$ R 对称性、Lorentz 对称性以及味对称性的 't Hooft 反常联系起来。反过来，'t Hooft 反常可以由库仑分支 EFT 通过反常匹配 [75] 确定。的确，SCFT 的（共形与味）中心荷会收到库仑分支上自由无质量场的贡献，但相关反常中还有 Wess--Zumino（WZ）贡献。概略地写为
\begin{equation}
  (a,c,k_G)_{\mathrm{SCFT}} = (a,c,k_G)^{\mathrm{free}}_{\mathrm{CB}} + \text{WZ 贡献}.
  \tag{3.42}
\end{equation}

% [Tutor Note: (3.42) 中带下标 CB、上标 free 的量指库仑分支上自由无质量场（$\mathrm{U}(1)$ 矢量多重态与带电超多重态）对 $a$、$c$、$k_G$ 的贡献；WZ 项是超出自由场部分的那些反常贡献。下一句（p.47 首句）指出，这些 WZ 贡献恰好重现了 (2.50) 中的 $R(A)$、$R(B)$。]
% [Tutor Note: 反常匹配之所以可行，是因为 't Hooft 反常沿整个 RG 流不变：可以在库仑分支上的弱耦合（自由场）一端把它算出来，再加上 WZ 项便得到 SCFT 的值——这就是 (3.42) 拆成「自由场部分 + WZ 部分」的原因。]

% [原文 p.47]
其中 WZ 贡献恰好重现了 (2.50) 中的 $R(A)$、$R(B)$ [75,145]。

\SEC{4}{五维膜给出的 Class S 构造}

本节中，我们探讨如何由 6d 理论的维数约化构造 4d $\cN=2$ 超对称理论。\footnote{关于 6d SCFT 的详细讨论超出了本综述的范围。读者可参阅 [181,182]，作为进入这一庞大主题文献的入口。}特别地，我们关注 Class S 理论（其中 “S” 代表 “Six”）。Class S 理论的出发点是 6d $\cN=(2,0)$ SCFT $\cT_{\fg}$，它由一个单链（simply-laced）李代数 $\fg$ 标记。然后我们把该理论紧化在黎曼面 $C$ 上，$C$ 称为 UV 曲线。这样的曲线上可以有穿刺点，对这些穿刺点我们需要为某些场指定边界条件，本节正文将对此作出说明。李代数、UV 曲线以及来自穿刺点 $p$ 的数据合在一起，定义一个 4d $\cN=2$ 的 Class S 理论，记作 $\cT(\fg,C,p)$。

% [Tutor Note: 原文的 “Class $\mathcal{S}$” 用花体 $\mathcal{S}$（正文与标题皆同），字面取 “Six” 之意——它来自 6d $\cN=(2,0)$ 理论在黎曼面 $C$ 上的紧化。翻译按术语表保留 “Class S” 英文；注意 $\cT_{\fg}$ 的 $\fg$ 是哥特体（fraktur），表示一般单链李代数，$C$ 是紧化用的黎曼面，$p$ 是穿刺点数据。]

我们将先在 4.1 节回顾 6d 的 $\cN=(2,0)$ 理论。4.2 节描述如何得到阿贝尔型 6d $\cN=(2,0)$ 理论的 Seiberg--Witten 解，并将在 4.3 节引入 UV 曲线 $C$ 与精确的 Class S 构造。4.4 节讨论通过求解经典 Hitchin 系统得到的量子预势，4.5 节引入 $C$ 上穿刺点的概念。4.6 节留给例子，4.7 节快速回顾 AGT 对应。最后，4.8 节简要提及进一步的推广。特别地，虽然前几节我们关注 $A_{N-1}$ 型理论，但在 4.8.1 节我们列出推广到其他结构群与 UV 曲线的可能途径；4.8.2 节引入更一般的穿刺点；最后在 4.8.3 节我们描述用理论李代数 $\fg$ 的外自同构进行扭曲（twisting）的步骤。

\SUB{4.1}{$\cN=(2,0)$ 理论的简要回顾}

6d $\cN=(2,0)$ SCFT $\cT_{\fg}$ 是可能存在的最大的超共形理论。一方面，6 是超对称与共形对称能够相容的最大维数 [97]。另一方面，它具有 6d 中可能的最大超共形代数，其实形式（real form）为 $\mathfrak{osp}(8^*|4)$ [183,184]。这个超代数包含 $\fso^*(8)\oplus\mathfrak{usp}(4)_R=\fso(2,6)\oplus\fso(5)_R$，它们分别是共形对称性与 R 对称性的玻色子代数。这样一些自由的 $\cN=(2,0)$ 理论由若干阿贝尔型张量多重态组成，每个多重态含一个实的自对偶 2-形式规范场 $B_{\alpha\beta}$、旋量 $\lambda$ 以及五个标量 $\Phi_I$。

% [Tutor Note: $\fso^*(8)\cong\fso(2,6)$ 是 6d 共形代数，$\mathfrak{usp}(4)_R\cong\fso(5)_R$ 是 R 对称性——后者正是五个标量 $\Phi_I$ 按 $\mathbf{5}$ 变换的来源。字体约定：本译本统一用哥特体（fraktur）表示李代数（$\mathfrak{osp}$、$\fso$、$\mathfrak{usp}$、$\fu$），大写正体表示群（$\mathrm{U}(1)$、$\SU(2)$）；原 PDF 此处印为粗体，仅为底本排印差异。]

% [原文 p.48]
相对于 R 对称性，它们分别按 $\fso(5)$ 的单态、$\mathbf{4}$ 与 $\mathbf{5}$ 不可约表示变换。

很长一段时间里，人们相信 6d 中不可能存在相互作用场论，因为拉格朗日量中的相互作用项在 $d\ge 5$ 维时不可重整。然而，人们已经认识到，相互作用的 6d $\cN=(2,0)$ 理论可以从弦/M 理论“工程”出来，例如在 M 理论中考虑一摞 M5 膜并去掉解耦的质心自由度。特别地，$N$ 个重合的 M5 膜以这种方式实现一个 $\cT_{\fsu(N)}$ SCFT，它没有拉格朗日量描述。

构造相互作用 6d $\cN=(2,0)$ SCFT 的另一种可能方式，是把 IIB 型弦论放在奇异几何 $\R^{1,5}\times\C^2/\Gamma$ 上，其中 $\Gamma\subset\SU(2)$ 是按 ADE 分类给出的 $\SU(2)$ 的有限子群 [185]。这一构造把前述 $\cT_{\fsu(N)}$ SCFT 作为特例 $\Gamma=\Z_N$ 包含在内。它也便于研究这些理论的模空间：我们可以通过把 $\C^2/\Gamma$ 原点处的奇点吹胀成一族有限大小的 2-闭链来沿模空间移动，而吹胀后的流形是一个超 Kähler ALE 流形。这种 2-闭链的数目等于代数 $\fg$ 的秩；对每个这样的例外 2-闭链，可以把上述五个标量 $\Phi_I$ 的 VEV 与 NS-NS 2-形式场、R-R 2-形式场以及三重辛形式的积分联系起来。于是 $\cT_{\fg}$ 的真空由这些 VEV 参数化，给出
\begin{equation}
  \cM_{\fg} = \R^{5r_{\fg}}/\cW_{\fg}\,,
  \tag{4.1}
\end{equation}
其中 $r_{\fg}$ 是 $\fg$ 的秩，$\cW_{\fg}$ 是它的 Weyl 群。

% [Tutor Note: (4.1) 的指数 $5r_{\fg}$ 来自每个例外 2-闭链上的 5 个标量与 $B$ 场积分（每个 2-闭链给出 5 个实数 VEV），而分母 $\cW_{\fg}$ 是 Weyl 群：沿模空间移动时，$\fg$ 的 Weyl 反射对应规范等价的重新参数化。$\cM_{\fg}$、$\cW_{\fg}$ 为花体，$\fg$ 为哥特体。]

这类理论的一个有用描述，是把它们紧化在半径为 $R$ 的圆 $S^1_R$ 上。阿贝尔型 6d 理论中的场 $(B,\lambda,\Phi)$ 约化为 $(A,\lambda,\Phi)$，其中 $A$ 是来自 $B$ 场约化的 5d 1-形式规范场，而 $\lambda$ 与 $\Phi$ 仍分别为费米子与标量，但现在已是 5d 的。所得理论是 5d $\cN=2$ 的阿贝尔型 SYM。更一般地，由于具有极大超对称，人们预期一般相互作用 6d $\cN=(2,0)$ SCFT 在 $S^1_R$ 上的紧化，在 Kaluza--Klein（KK）标度以下由 5d $\cN=2$ 非阿贝尔型 SYM 描述，其规范耦合 $g^2_{\mathrm{YM}}\sim R$。事实上，通过匹配 5d 库仑分支上的某些 BPS 态与 6d 张量分支上的 BPS 态，可以重新得到 6d $\cN=(2,0)$ SCFT 的 ADE 分类 [186]。重要的是，5d SYM 会通过其瞬子粒子悄悄地记住 6d 的圆：这些瞬子粒子在拓扑流 $J=\star\tr(F\wedge F)$ 下带电，质量 $m_I\sim\frac{1}{g^2_{\mathrm{YM}}}$。它们自然被等同于 $S^1_R$ 紧化的 KK 模 [187,188]。这一特征使得人们能够通过追踪 KK 塔

% [原文 p.49]
（例如 [189--191] 中的 6d 超共形指标），从 5d SYM 确定 6d $(2,0)$ SCFT 中的受保护可观测量。\footnote{顺便说一句，从这里可以看出为什么 $(2,0)$ 理论没有明显的拉格朗日量描述：对 6d 作用量做朴素的维数约化，得到的 5d 作用量正比于 $R$，而不是反比于 $R$。关于反对拉格朗日量存在的更多理由，见 [192,193]。}

% [Tutor Note: 注意 $g^2_{\mathrm{YM}}\sim R$ 是正比而非反比，这是 KK 约化的标志；瞬子粒子质量 $m_I\sim1/g^2_{\mathrm{YM}}\sim1/R$ 与 KK 塔的首个能隙一致，故它们被认定为 KK 模。脚注 34 正是用这一点解释“为何不存在明显的拉格朗日量”：朴素约化直接给出正比于 $R$ 的 5d 作用量。]

\SUB{4.2}{阿贝尔 $\cN=(2,0)$ 理论给出的 Seiberg--Witten 解}

% [Tutor Note: 实测 p.50（印刷 p.49）上 §4.2 的正文完整（四段 + 脚注 34），(4.2)(4.3) 在 p.51 且属 §4.3，并非 §4.2 的公式；BRIEF/TIMELINE 的「4.2 正文在 p.51」与实测不符，已回执报告 Lead。]

还可以再往前走一步：把该理论紧化在另一个圆上（一般地，我们将紧化在黎曼面上），并自然地预期在低能极限下得到一个 4d QFT。在进入 Class S 构造之前，让我们看看这与 Seiberg--Witten 的故事有什么关系。

回忆第 3 节：一般 4d $\cN=2$ SCFT 的 EFT 由其 SW 几何编码。聚焦于秩为 1 的理论，这一几何由一条纤维化在 $u$ 平面上的椭圆曲线构成（$u$ 平面由 EFT 中 4d 矢量多重态标量的 VEV 参数化）。

有效作用可以由上述曲线 $\Sigma_u$ 与 SW 微分 $\lambda$ 描述。$(\Sigma_u,\lambda)$ 这一对对象确定了预势，而预势包含长波极限下库仑分支上动力学的信息。这一几何结构在 6d 理论中有物理含义。我们可以考察约化在 SW 曲线 $\Sigma_u$ 上的 6d $\fu(1)$ $\cN=(2,0)$ 理论，由此得到 4d $\cN=2$ EFT。特别地，4d EFT 有一个规范场及其对偶 $(A_\mu,A^D_\mu)$，它们来自 6d $\cN=(2,0)$ 理论中自对偶 2-形式 $B$ 在 $H_1(\Sigma_u)$（及其 Hodge 对偶）的典范基上的约化。此外，由 3.3 节我们知道，BPS 粒子的质量由 Seiberg--Witten 微分 $\lambda$ 在 1-闭链上的积分给出。从 6d 的视角看，这类 BPS 粒子来自缠绕这些 1-闭链的 BPS 弦。

这一点可以从 M 理论的视角很好地理解：那里有一个单独的 M5 膜包裹在 SW 曲线上。BPS 弦来自某些超对称的 M2 膜，而这些 M2 膜（它们还沿另一个横向方向延展）的无穷小张力自然地给出 Seiberg--Witten 微分 $\lambda$。若 4d 理论的秩 $r>1$，以致在库仑分支的一般点上规范群破缺为 $\mathrm{U}(1)^r$，则唯一的差别是：这个单独 M5 膜所包裹的 SW 曲线不再是椭圆曲线，而是亏格等于秩的曲线 [28]。

至此，我们把 4d 库仑分支的有效物理提升为一个紧化在黎曼面上的 6d 自由理论。下一步将说明，借助 Class S 构造，这样的 6d 理论具有相互作用的 UV 完备化。

% [边界] 本块止于 p.50（印刷 p.49）页末。p.51 首行是 §4.3 标题与正文（属 tr-05）；(4.2)(4.3) 在 p.51 的 §4.3 中。
